\chapter{Growth Hormone (GH)}

\section*{What Is This Test?}
Growth hormone (GH, also called somatotropin) is a 191-amino-acid single-chain polypeptide hormone produced by somatotrope cells of the anterior pituitary gland. GH secretion is regulated by two hypothalamic hormones: GH-releasing hormone (GHRH), which stimulates GH synthesis and secretion, and somatostatin (also called somatotropin release-inhibiting hormone), which inhibits GH release. Ghrelin, a peptide produced by the stomach, also stimulates GH secretion. GH is secreted in a pulsatile pattern, with 6--10 pulses per 24 hours, the largest occurring during slow-wave (stage N3) sleep \cite{tietz2023textbook}. Between pulses, GH levels are undetectable (<0.1 ng/mL) in normal individuals. GH exerts its growth-promoting and metabolic effects both directly (via GH receptors) and indirectly through stimulating hepatic and peripheral production of insulin-like growth factor 1 (IGF-1). GH has anabolic effects: stimulates linear growth in children, increases protein synthesis, promotes lipolysis, and antagonizes insulin action (raising blood glucose). The GH test measures the concentration of GH in serum. However, because of pulsatile secretion and the wide range between pulses (undetectable to 30--50 ng/mL), a single random GH measurement has essentially no diagnostic value. GH must be measured in the context of dynamic testing: suppression testing (oral glucose tolerance test) for acromegaly and stimulation testing (insulin tolerance test, GHRH-arginine, glucagon, clonidine, or L-dopa) for GH deficiency.

\section*{Alternative Names}
Somatotropin; Human Growth Hormone (hGH); GH, Serum; GH Suppression Test; Oral Glucose Tolerance Test for GH (OGTT-GH); GH Stimulation Test; Insulin Tolerance Test (ITT); GHRH-Arginine Stimulation Test; Glucagon Stimulation Test

\section*{Principle}
Serum GH is measured using two-site sandwich immunometric (chemiluminescent) immunoassays on automated analyzers. Two monoclonal antibodies against different epitopes on the GH molecule (typically recognizing the 22 kDa monomeric GH isoform, which represents 75--85\% of circulating GH) are used. One antibody (capture) is bound to a solid phase, and the second antibody (detection) is labeled with a chemiluminescent tag (ruthenium on Roche, acridinium ester on Abbott/Siemens). After incubation and washing, the signal is proportional to GH concentration. GH assays vary significantly across manufacturers (inter-assay CV 10--25\%), largely because of differing antibody specificities for GH isoforms (22 kDa, 20 kDa, and other variants), GH-binding protein interference, and differences in calibration against the recombinant WHO International Standard (IS 98/574). The 2011 consensus statement recommended the use of 22 kDa-specific GH assays calibrated to the recombinant international standard \cite{clemmons2011consensus}. Functional sensitivity of high-sensitivity GH assays is <0.05 ng/mL, which is essential for detecting nadir GH during an oral glucose tolerance test for acromegaly (nadir GH <1 ng/mL is considered normal suppression; <0.4 ng/mL by ultrasensitive assay). LC-MS/MS methods for GH have been developed but are not in routine clinical use.

\section*{History}
GH was first isolated from bovine pituitary glands by Herbert Evans and colleagues in the 1920s--1930s, demonstrating that pituitary extracts stimulated growth in rats. Human GH was first purified from cadaveric pituitary glands by Choh Hao Li and Harold Papkoff in the 1950s. Human cadaveric GH was used therapeutically to treat GH-deficient children beginning in the late 1950s but was withdrawn in 1985 after cases of Creutzfeldt-Jakob disease were linked to contaminated pituitary preparations (over 200 cases worldwide). Recombinant human GH (somatropin) was developed by Genentech using recombinant DNA technology and approved by the FDA in 1985, eliminating the infectious risk. The radioimmunoassay for GH was developed by Utiger and Odell in 1962. The clinical syndrome of acromegaly was first described by Pierre Marie in 1886. The use of the oral glucose tolerance test to diagnose acromegaly (failure of GH to suppress below <1 ng/mL) was established in the 1970s. The 2014 Endocrine Society acromegaly guidelines \cite{katznelson2014acromegaly} and the 2016 Pediatric Endocrine Society guidelines for GH stimulation testing in children \cite{grimberg2016guidelines} standardized the diagnostic approach.

\section*{Why This Test Is Ordered}
\begin{itemize}
  \item Diagnosis of acromegaly (GH excess in adults) and gigantism (GH excess before epiphyseal closure in children) --- oral glucose tolerance test (75 g glucose) with GH measurement at 0, 30, 60, 90, and 120 minutes; failure of GH to suppress below <1 ng/mL confirms acromegaly (or <0.4 ng/mL with ultrasensitive assays)
  \item Evaluation of suspected GH deficiency in children with short stature (height >2 standard deviations below the mean, height velocity <25th percentile, or height below mid-parental target) --- GH stimulation testing using 2 different provocative agents
  \item Diagnosis of adult GH deficiency in patients with known hypothalamic-pituitary disease, previous pituitary surgery/radiation, traumatic brain injury, or history of childhood-onset GH deficiency transitioning to adulthood --- insulin tolerance test (ITT) preferred, or GHRH-arginine, or glucagon stimulation
  \item Monitoring disease activity after treatment of acromegaly (transsphenoidal surgery, somatostatin analogues, pegvisomant, radiation) --- nadir GH after OGTT <1 ng/mL or random GH <1 ng/mL AND normal age-adjusted IGF-1 indicates remission
  \item Assessing adequacy of GH replacement therapy --- target IGF-1 in the mid-normal range for age, with GH levels maintained within physiologic range
  \item Investigation of suspected ectopic GHRH or GH secretion causing acromegaly (rare: bronchial carcinoid, pancreatic islet cell tumors, hypothalamic hamartomas)
\end{itemize}

\section*{Primary vs Secondary Test}
Random serum GH is NOT a diagnostic test and has no clinical value as a standalone measurement. GH measurement as part of dynamic testing (suppression or stimulation) is a secondary test used to confirm suspected GH excess or deficiency. For acromegaly, the primary screening test is IGF-1, and GH is measured during the OGTT for confirmation. For GH deficiency, stimulation testing with GH measurement is required for definitive diagnosis, but the decision to test is based on clinical evaluation, auxological data, and often low IGF-1/IGFBP-3 levels \cite{molitch2011evaluation}.

\section*{How This Test Is Performed}
For a random GH: a blood sample is drawn from a vein into a serum separator tube (gold-top). For the OGTT for acromegaly: the patient drinks 75 g of glucose solution. Serum GH and glucose are measured at baseline (0 minutes) and at 30, 60, 90, and 120 minutes after ingestion. In normal individuals, hyperglycemia suppresses GH to <1 ng/mL (or <0.4 ng/mL with ultrasensitive assays). In acromegaly, GH fails to suppress (nadir >1 ng/mL) and may paradoxically rise. For insulin tolerance test (ITT) for GH deficiency: IV insulin 0.1 units/kg is administered to induce hypoglycemia (glucose <40 mg/dL or <2.2 mmol/L). GH is measured at baseline and at 15, 30, 45, 60, 90, and 120 minutes. Peak GH <5 ng/mL in adults or <10 ng/mL in children indicates GHD. This test requires continuous physician supervision, IV access, and immediate glucose rescue capability. For GHRH-arginine stimulation test: GHRH 1 mcg/kg IV bolus followed by arginine 0.5 g/kg IV (max 30 g) over 30 minutes. GH measured at baseline and at 15, 30, 45, 60, 90, and 120 minutes. For glucagon stimulation test: glucagon 1 mg IM (or 1.5 mg if >90 kg), GH measured at baseline and at 30, 60, 90, 120, 150, 180, 210, and 240 minutes.

\section*{What Preparation Is Needed}
The patient must fast for 10--12 hours overnight before any dynamic GH testing. Random GH draws do not require fasting but are not clinically useful. The patient should avoid strenuous exercise for 24 hours before testing because exercise is a potent GH secretagogue. For the OGTT, the glucose load (75 g in 300 mL) must be consumed within 5 minutes. For the ITT, the patient must be monitored continuously by a physician with IV access, and 50\% dextrose (D50W) must be immediately available for rescue from severe hypoglycemia. Contraindications to ITT include: coronary artery disease, history of seizure disorder, pregnancy, and age >65 years. Estrogen therapy (oral contraceptives, HRT) increases GH binding protein and alters GH response to stimulation testing --- oral estrogen should be discontinued for 6--8 weeks before GH stimulation testing in postmenopausal women being evaluated for adult GHD. Obesity markedly blunts GH responses to all stimulation tests; BMI-adjusted cut-offs for peak GH should be used (obesity reduces GH peak by 50--70\%). Male patients should be adequately androgenized (testosterone replacement if hypogonadal) before GH stimulation testing.

\section*{Contraindications and Precautions}
The most important pre-analytical consideration is that random serum GH is almost never interpretable. GH secretion is pulsatile and episodic, with undetectable levels between pulses. A single random GH level of 10 ng/mL could represent a physiologic pulse or acromegaly; a level of <0.1 ng/mL could represent the inter-pulse nadir or GH deficiency. Stress (venipuncture, hospitalization, anxiety, acute illness) is a powerful GH secretagogue and markedly increases GH. Exercise, protein meal, hypoglycemia, and sleep all stimulate GH. Drugs affecting GH: L-dopa, clonidine, arginine, GHRH, and glucagon stimulate GH; somatostatin analogues (octreotide, lanreotide) and the GH receptor antagonist pegvisomant lower GH; dopamine agonists (bromocriptine, cabergoline) paradoxically suppress GH in acromegaly; insulin suppresses GH via hyperglycemia; glucocorticoids suppress GH; estrogen increases GH binding protein; obesity profoundly blunts GH secretion (BMI inversely correlated with peak GH on stimulation). Age: GH secretion declines progressively with age after the third decade (somatopause), so normative values for GH stimulation peaks are age-adjusted. Pregnancy: placental GH (GH-V) replaces pituitary GH (GH-N) by the second trimester, and GH assays must recognize that GH-V cross-reacts variably. Manufacturer variability: GH assays vary 2--3 fold in reported GH concentrations for the same specimen because of differing isoform specificities and calibrator standardization. The same assay and laboratory MUST be used for serial monitoring.

\section*{Risks}
From venipuncture: slight pain, bruising. For OGTT: nausea, bloating, vomiting (common with 75 g glucose load). For ITT: severe hypoglycemia (required to provoke GH release), with symptoms of neuroglycopenia (confusion, seizure, loss of consciousness), sweating, tachycardia, tremor. ITT is contraindicated in patients with coronary artery disease (can provoke myocardial ischemia), epilepsy, pregnancy, and age >65. ITT-related mortality is approximately 0.01\% (1 in 10,000), and ITT should ONLY be performed in experienced endocrine testing units with continuous physician supervision. For glucagon stimulation: nausea and vomiting (common, 30--50\%), mild hyperglycemia. For GHRH-arginine: facial flushing, transient hypotension from arginine. For clonidine: hypotension and sedation (monitor BP and heart rate for 3 hours).

\section*{What Happens After the Test}
GH results from dynamic testing are available within 1--3 days (GH measurement is routine on automated platforms; dynamic testing requires coordination with the laboratory). For acromegaly: failure of GH to suppress below 1 ng/mL after 75 g glucose AND elevated age-adjusted IGF-1 confirms the diagnosis. Next step is pituitary MRI with gadolinium. A pituitary macroadenoma (>1 cm) is found in 75--80\% of acromegaly cases at presentation; microadenomas in the remainder \cite{melmed2006acromegaly}. For GH deficiency: peak GH <5 ng/mL on ITT (or <3 ng/mL with BMI >25, <1 ng/mL with BMI >30) or <4.1 ng/mL on GHRH-arginine (with BMI-adjusted cut-offs) confirms adult GHD. Next step: pituitary MRI to identify cause (empty sella, pituitary adenoma, pituitary hypoplasia), and evaluation of other pituitary hormones. GH replacement therapy (somatropin) can then be initiated at low doses (0.1--0.3 mg/day) and titrated based on IGF-1 response and clinical benefit.

\section*{Understanding Results}

\subsection*{Below Normal}
\begin{itemize}
  \item \textbf{Peak GH <5 ng/mL on ITT (adult criteria, with hypoglycemia <40 mg/dL confirmed):} Growth hormone deficiency (adult). Causes: pituitary tumors (non-functioning adenoma, craniopharyngioma, Rathke cleft cyst), pituitary surgery or radiation, traumatic brain injury, Sheehan syndrome, lymphocytic hypophysitis, empty sella syndrome, genetic causes (POU1F1 [Pit-1], PROP1 mutations), or idiopathic.
  \item \textbf{Peak GH <10 ng/mL on 2 stimulation tests (pediatric criteria):} Childhood GH deficiency. Associated with short stature (>2 SD below mean for age/sex), decreased growth velocity (<25th percentile), delayed bone age, increased subcutaneous fat, mid-face hypoplasia, and micropenis in males. Treatment: recombinant human GH (somatropin) 0.025--0.05 mg/kg/day subcutaneously.
  \item \textbf{Undetectable or very low GH (<0.1 ng/mL) after pituitary surgery for acromegaly:} Surgical remission --- indicates the somatotroph adenoma has been completely removed. If accompanied by normal age-adjusted IGF-1, this is biochemical cure.
  \item \textbf{Suppressed GH (<1 ng/mL nadir) on OGTT:} Normal --- effectively excludes acromegaly in patients without confounding conditions (see above normal section for exceptions).
  \item \textbf{Physiologic low GH in obesity:} Obese individuals have markedly blunted GH responses to all stimulation tests. BMI-adjusted cut-offs for peak GH must be used (BMI 25--30: peak GH <3 ng/mL; BMI >30: peak GH <1 ng/mL).
\end{itemize}

\subsection*{Above Normal}
\begin{itemize}
  \item \textbf{Failure of GH to suppress below 1 ng/mL after 75 g OGTT:} Diagnostic of acromegaly (with elevated age-adjusted IGF-1). In normal individuals, hyperglycemia (glucose >140 mg/dL) suppresses GH to <1 ng/mL within 60--120 minutes. In acromegaly, GH paradoxically remains elevated or rises. Some patients show partial suppression (nadir 1--5 ng/mL); virtually all acromegaly patients have nadir >1 ng/mL on ultrasensitive assay (<0.4 ng/mL).
  \item \textbf{Acromegaly (GH-secreting pituitary adenoma, somatotroph adenoma):} >99\% of acromegaly is caused by a pituitary adenoma \cite{melmed2009acromegaly}. GH is elevated (often >5 ng/mL randomly, but can be lower in microadenomas). IGF-1 is elevated above age-adjusted upper limit of normal. Mixed GH/prolactin-secreting adenomas (mammosomatotroph) occur in 20--30\%. Excess GH causes: soft tissue swelling (coarsening of facial features, enlargement of hands and feet --- increase in ring and shoe size), macroglossia, prognathism (jaw protrusion), frontal bossing, oily skin, skin tags, carpal tunnel syndrome, osteoarthritis, obstructive sleep apnea, hypertension, cardiomyopathy, diabetes mellitus, and increased risk of colorectal neoplasia (colonoscopy recommended at diagnosis and every 5--10 years). Mortality is 2--4 fold increased compared to general population, normalizing after biochemical remission.
  \item \textbf{Ectopic GHRH secretion:} Rare (<1\% of acromegaly). Bronchial or pancreatic carcinoid tumors secreting GHRH stimulate pituitary somatotrope hyperplasia (not adenoma). Serum GHRH levels are elevated. CT/MRI chest and abdomen should be performed.
  \item \textbf{Gigantism:} GH excess before closure of epiphyseal growth plates (before puberty) causes excessive linear growth (height >3 SD above mean, may reach 7--8 feet). Causes: pituitary gigantism (somatotroph adenoma or hyperplasia), McCune-Albright syndrome (GNAS mutation causing constitutive Gs-alpha activation), Carney complex (PRKAR1A mutation), and X-linked acrogigantism (Xq26.3 microduplication including GPR101 gene, causing early-onset gigantism).
  \item \textbf{Artifactual GH elevation from stress:} Stress-induced GH release during venipuncture, acute illness, surgery, or trauma can produce GH >10--50 ng/mL. This is a false-positive. Dynamic testing with glucose suppression is required to exclude acromegaly.
  \item \textbf{Pregnancy (second and third trimester):} Placental GH (GH-V) replaces pituitary GH (GH-N). Most GH assays cross-react with GH-V, so GH levels appear elevated. This is physiologic and reverts to normal after delivery.
  \item \textbf{Uncontrolled diabetes mellitus:} GH can be mildly elevated because insulin deficiency reduces IGF-1 production in liver, reducing negative feedback and increasing GH secretion. GH elevation resolves with improved glycemic control.
  \item \textbf{Malnutrition and anorexia nervosa:} GH is elevated with low IGF-1 (GH resistance state due to malnutrition). GH decreases with nutritional rehabilitation.
\end{itemize}

\section*{Normal Results}
\begin{itemize}
  \item \textbf{Random serum GH (resting, fasting):} <0.1--0.5 ng/mL (undetectable between pulses) --- a single random GH is not clinically interpretable
  \item \textbf{GH nadir during OGTT (75 g glucose):} <1.0 ng/mL (<0.4 ng/mL with ultrasensitive assay) --- excludes acromegaly
  \item \textbf{GH in acromegaly on OGTT:} Nadir >1 ng/mL (failure to suppress)
  \item \textbf{GH stimulation, insulin tolerance test (ITT):} Peak GH >5 ng/mL (adult); >10 ng/mL (child) --- adequate GH reserve
  \item \textbf{GH in adult GHD (ITT):} Peak GH <5 ng/mL (confirms severe GHD); BMI-adjusted: <3 ng/mL (BMI 25--30), <1 ng/mL (BMI >30)
  \item \textbf{GH stimulation, GHRH-arginine:} Peak >4.1 ng/mL (normal); <4.1 (GHD, BMI-adjusted cut-offs apply)
  \item \textbf{GH stimulation, glucagon:} Peak >3 ng/mL (normal); <3 (GHD)
  \item \textbf{Post-treatment acromegaly (remission):} Random GH <1 ng/mL AND normal age-adjusted IGF-1
  \item \textbf{Note:} GH assays are not harmonized across manufacturers. Cut-off values are assay-specific. Always use the same assay for serial monitoring, and interpret using the assay-specific reference interval.
\end{itemize}

\section*{Follow-up Tests}
\begin{itemize}
  \item \textbf{IGF-1 (insulin-like growth factor 1):} The primary screening test for acromegaly. IGF-1 integrates GH secretion over 24 hours, is stable (no pulsatility), and should always be measured with GH. Elevated age- and sex-adjusted IGF-1 is the hallmark of acromegaly. IGF-1 is also used to monitor treatment response.
  \item \textbf{Oral glucose tolerance test with GH (OGTT-GH):} Gold standard confirmatory test for acromegaly. Nadir GH >1 ng/mL confirms the diagnosis. The glucose load must achieve hyperglycemia (>140 mg/dL) to properly suppress GH.
  \item \textbf{Pituitary MRI with gadolinium contrast:} Thin-section (1--3 mm), T1-weighted images before and after contrast in coronal and sagittal planes. Detects pituitary adenoma (somatotroph macroadenoma >1 cm in 75--80\%). Evaluates tumor size, extension (optic chiasm compression, cavernous sinus invasion), and resectability.
  \item \textbf{Insulin tolerance test (ITT):} Gold standard for diagnosing adult and childhood GH deficiency. Hypoglycemia (glucose <40 mg/dL) is a potent GH secretagogue. Contraindicated in seizure disorders, CAD, pregnancy, age >65.
  \item \textbf{GHRH-arginine stimulation test:} Alternative to ITT when ITT is contraindicated. Less validated in hypothalamic causes of GHD (GHRH directly stimulates pituitary, may give false-normal result in hypothalamic GHD).
  \item \textbf{Glucagon stimulation test:} Alternative stimulation test for GH deficiency. Particularly useful when ITT and GHRH-arginine are contraindicated or unavailable. Requires 3--4 hours with multiple blood draws.
  \item \textbf{IGFBP-3 (IGF-binding protein 3):} Less sensitive than IGF-1 but can support the diagnosis of GH deficiency in children. Low IGF-1 and low IGFBP-3 suggest GHD; low IGF-1 with normal IGFBP-3 suggests malnutrition or other causes.
  \item \textbf{Bone age X-ray (left hand and wrist):} For childhood growth evaluation. Delayed bone age (>2 SD below chronological age) is characteristic of GH deficiency.
  \item \textbf{Colonoscopy:} For all patients with acromegaly at diagnosis and periodically thereafter, because of 2--3 fold increased risk of colorectal adenomas and carcinoma. Screen at diagnosis and every 5--10 years depending on findings.
  \item \textbf{Echocardiogram:} For acromegaly patients to evaluate for GH excess-related cardiomyopathy (biventricular hypertrophy, diastolic dysfunction).
  \item \textbf{Sleep study (polysomnography):} For acromegaly patients with symptoms of obstructive sleep apnea (macroglossia and soft tissue swelling narrow the upper airway).
\end{itemize}

\section*{References}
\bibliographystyle{unsrtnat}
\bibliography{references}

\clearpage
\section*{Regulatory Status and Authorization Date}
Regulatory status applies to the specific assay, instrument, manufacturer, and intended use rather than to the test name alone. No single approval date can be assigned to this test category. For a commercial assay, verify the current product in the relevant regulator's database: the U.S. Food and Drug Administration (FDA) clearance, approval, or authorization; India's Central Drugs Standard Control Organisation (CDSCO) licence or permission; the European Union In Vitro Diagnostic Regulation (EU IVDR) CE certificate issued through the applicable conformity-assessment route; or the United Kingdom Medicines and Healthcare products Regulatory Agency (MHRA) registration and UKCA/CE status. Laboratory-developed tests may not have an individual FDA, CDSCO, EU IVDR, or MHRA approval date.
\clearpage
